\documentclass[10pt,a4paper]{amsart}
\usepackage[margin=19mm]{geometry}
\usepackage{amsmath,amssymb,mathtools}
\usepackage[scr=rsfs,cal=euler]{mathalfa}
\usepackage{xfrac}
\usepackage[unicode,hidelinks,bookmarks=false]{hyperref}
\newcommand{\ie}{i.e.\ }
\newcommand{\cf}{cf.\ }
\newcommand{\resp}{respectively\ }
\newcommand{\Subsection}{\subsection}
\providecommand{\nobreakdash}{\nobreak\hbox{-}}
\setlength{\emergencystretch}{2em}
\multlinegap0pt
\usepackage{bm}
\usepackage{bbm}
\usepackage{dsfont}
\usepackage{xspace}
\usepackage{cancel}
\usepackage{mathtools}
\usepackage{amsthm}
\makeatletter
\@ifundefined{theorem}{\newtheorem{theorem}{Theorem}}{}
\makeatother
\title[Coherent Local Explanations for Mathematical Optimization]{Coherent Local Explanations for Mathematical Optimization}
\author{Daan Otto, Jannis Kurtz, S. Ilker Birbil}
\date{Source: arXiv:2502.04840v2 (CC BY 4.0)}
\begin{document}
\maketitle

\noindent\emph{Source and licence.} Coherent Local Explanations for Mathematical Optimization. Version arXiv:2502.04840v2 (CC BY 4.0), distributed under the Creative Commons Attribution 4.0 International licence, as recorded in the arXiv licence metadata for this exact version and shown on the abstract page https://arxiv.org/abs/2502.04840v2. Full rights notice: licenses/arxiv-2502.04840-CC-BY-4.0.txt.\par\smallskip

\noindent\textit{Editorial scope.} Counted unit: the Theorem whose source label is \texttt{thm:provably\_coherent} in the article (source0.tex lines 754--763, source label \texttt{thm:provably\_coherent}), delivered together with the article's own complete original proof of it (source0.tex lines 764--785, one \texttt{proof} environment of 22 lines). No mathematics is written, repaired or re-derived: statement and proof are the article's own text line for line. Required context retained uncounted: eq:coherence1 (lines 208--211). Every definition, display and label of those lines is retained unchanged. Omitted from this package and not redistributed: 1--207, 212--753, 786--871. Nothing inside a retained excerpt is omitted or altered: every retained body is delivered exactly as the article's lines are written, so no omission receipt of any kind accompanies this package. Citations: the retained excerpts carry no citation command at all, so no bibliography entry of the article is transcribed and no reference is redistributed. Reference adaptation: none was needed and none was made; every reference inside the delivered unit resolves inside it, so no unresolved reference or citation is produced. Printed slips kept literal: no printed word, symbol, hypothesis or number of the counted statement or of its proof was altered. Layout and class: the article's own class is \texttt{article} with packages including Arxiv, inputenc, fontenc, hyperref, url, booktabs, amsfonts, nicefrac, microtype, xcolor, graphicx, tikz, subfigure, multirow; the wrapper is the lane's pinned \texttt{amsart} editorial wrapper at 10pt on A4, which restates the theorem-like environments the retained text needs and adds the article's own macro definitions that the retained text needs (none), with extra wrapper packages (bm, bbm, dsfont, xspace, cancel, mathtools). The delivered numbering is the unit's own. Assets: no retained span contains a figure environment, a graphics-inclusion command, a PDF-page inclusion, a table or a TikZ picture; no image file is copied into this package, the package declares no asset dependency and no image file is redistributed with it. Licence: version arXiv:2502.04840v2 of this article is distributed under the Creative Commons Attribution 4.0 International licence, as recorded in the arXiv licence metadata for this exact version and shown on the abstract page https://arxiv.org/abs/2502.04840v2. A full rights notice is delivered at licenses/arxiv-2502.04840-CC-BY-4.0.txt. Characters: every retained excerpt is pure ASCII, so the delivered engine, XeLaTeX with its default Latin Modern text font, reports no missing character.
\begin{gather}
    f(g_{\bm{x}}(\bm{\theta}); \bm{\theta}) = g_{f}(\bm{\theta}), \label{eq:coherence1}\\
    g_{\bm{x}}(\bm{\theta}) \in \mathbb X(\bm{\theta}).\label{eq:coherence2}
\end{gather}

\begin{theorem}\label{thm:provably_coherent}
The minimizers of the weighted least-square problems
\[
\min_{\bm{\beta}_{f}} \ \sum_{i=0}^{N} w^i\| f(\bm{x}^i;\bm{\theta}^i) - \bm{\beta}_{f}^\intercal\bm{\theta}^i \|^2 \]
and 
\[
\min_{\bm{\beta}_{x_j}} \ \sum_{i=0}^{N} w^i\| \bm{x}^i_j - \bm{\beta}_{x_j}^\intercal\bm{\theta}^i \|^2 \quad j=1,\ldots ,p
\]
fulfill the coherence condition in \eqref{eq:coherence1}.
\end{theorem}

\begin{proof}
Since the objective function $\bm{\hat c}^\intercal \bm{x}$ is fixed and linear we have
\[
f(\bm{x}^i;\bm{\theta}^i) = \sum_{j=1}^{p} \hat c_j \bm{x}^i_j.
\]
The weighted least-squares problem has the unique optimal solution
\[
\bm{\beta}_{x_j}^* = (\bm{\Theta}^\intercal \bm{W}\bm{\Theta})^{-1} \bm{\Theta}^\intercal \bm{W} \bm{y^j} \quad j=1,\ldots ,p
\]
and 
\[
\bm{\beta}_{f}^* = (\bm{\Theta}^\intercal \bm{W}\bm{\Theta})^{-1} \bm{\Theta}^\intercal \bm{W} \bm{y^f},
\]
where $\bm{\Theta}$ is the matrix whose $i$-th row is the vector $\bm{\theta}^i$, $\bm{W}$ is the matrix with weight $w^i$ on the diagonal and zeroes elsewhere, $\bm{y^j}$ is the vector where the $i$-th entry is the value $x_j^i$ and $\bm{y^f}$ is the vector where the $i$-th entry is the value $f(\bm{x}^i;\bm{\theta}^i)$. We assume here that $(\bm{\Theta}^\intercal \bm{W}\bm{\Theta})$ is invertible. Then for any new parameter vector $\bm{\theta}$ the predicted optimal value of our model is
\begin{align*}
\bm{\theta}^\intercal\bm{\beta}_{f} & = \bm{\theta}^\intercal (\bm{\Theta}^\intercal \bm{W}\bm{\Theta})^{-1} \bm{\Theta}^\intercal \bm{W} \bm{y^f}  \\
& = \bm{\theta}^\intercal (\bm{\Theta}^\intercal \bm{W} \bm{\Theta})^{-1} \bm{\Theta}^\intercal \bm{W} \left( \sum_{j=1}^{p} \hat c_j \bm{y^{j}}\right) \\
& = \sum_{j=1}^{p} \hat c_j \bm{\theta}^\intercal (\bm{\Theta}^\intercal \bm{W}\bm{\Theta})^{-1} \bm{\Theta}^\intercal \bm{W} \bm{y^{j}} \\
& = \sum_{j=1}^{p} \hat c_j \bm{\theta}^\intercal\bm{\beta}_{x_j},
\end{align*}
which means that the predictors are coherent regarding condition \eqref{eq:coherence1}. 
\end{proof}

\end{document}
